\scanpage{10}
% source lines 1-4, continuation
\unsure{je} \unsure{considère|considérer} un $a\in K^*$ tel que
\[
\left\{\begin{aligned}
a&\equiv1\ (\underline m)\quad\text{aux places }\deleted{\unsure{\mathfrak p}}\ \unsure{\mathfrak q|v}\mid\underline m,\ \mathfrak q\neq\mathfrak p\\
a&\equiv x\quad\underline m\quad\text{en }\mathfrak p
\end{aligned}\right.
\]
[\deleted{\unsure{la définition}} Si $a'$ en est un autre, alors $aa'^{-1}\deleted{\unsure{\in K}}\equiv1\ (\underline m)$]

On prend alors $\operatorname{div}_{\complement\underline m}(a)\bmod P_{\underline m}(K)$.

% source lines 5-12
\uline{\unsure{Ènoncé|Énoncé} des \unsure{théorémes|théorèmes} de Artin--Takagi}. Soit $L$ une extension \unsure{galoisienne} \inserted{finie} de $K$. \deleted{\unsure{Pour de}} Soit $S$ l'ens des places de $K$ où $L$ est ramifié. Pour $\mathfrak p\in\complement S$, on définit $\sigma(\unsure{\mathfrak P|\mathscr P})\in\operatorname{Gal}(L/K)=G$ \inserted{(substitution de Frobenius)} dans le groupe de décomposition d'un $\mathfrak P$ au-dessus de $\mathfrak p$ (si on change $\unsure{\mathfrak p|\mathfrak P}$, $\sigma(\mathfrak P)$ ne change que par un aut. intérieur), et définit donc un $\unsure{=}$ élément de $G/G'$. \deleted{Alors} \inserted{On} prolonge par linéarité en un hom.
\[
{\unsure{D|\mathscr{D}}}_{\underline m}(K)\longrightarrow G/G'\qquad(\text{hom. de réciprocité})
\]

% source lines 13-17, circled 1 degree
\textcircled{$1^\circ$} \uline{\uuline{On peut trouver un module}} $\underline m$, de support $S$, t.q. l'hom de réciprocité s'annule sur $P_{\underline m}(K)$.

On peut donc le considérer comme un hom
\[
C_{\underline m}(K)\longrightarrow G/G',
\]
\inserted{donc $C(K)\longrightarrow G/G'$}
(cet dernier ne dépend plus du choix de $\underline m$]
\margin{[\unsure{n'a} de sens que dans $G/G'$ !]}

% source lines 18-23, circled 2 degree
\textcircled{$2^\circ$} \uline{Cet hom. est surjectif} [de façon plus précise, pour $\forall$ classe d'\unsure{élmts|élts} de $G$, il existe une \uline{infinité} de places $\mathfrak p$ pour \unsure{lesquels|lesquelles} $\sigma(\mathfrak p)$ soit dans cette classe]. Ainsi, $G/G'$ est un \uline{quotient} de $C_{\underline m}(K)$

% source line 24, continuing on scan 11
\textcircled{$3^\circ$} Le \uline{noyau de l'hom de réciprocité} est
% archive sheet number: 9
